\chapter{Metodologia}\label{chp:met}

\section{Visão Geral}

A estimação horária de emissões de \ac{CO$_2$} em usinas termelétricas configura-se como um problema inverso: dispõe-se de medições anuais agregadas, fornecidas pelos inventários do \ac{IEMA}, mas necessita-se do comportamento horário para aplicações em planejamento energético e despacho econômico-ambiental. Além disso, a relação entre geração horária e emissão é intrinsecamente não-linear \cite{gent1971minimum, talaq1994summary}.

Para enfrentar esses desafios, a metodologia proposta estrutura-se em duas etapas complementares. Na primeira etapa, resolve-se um problema de otimização não-linear para calibrar os coeficientes de uma equação paramétrica de emissão horária, utilizando como referência os totais anuais disponibilizados pelo \ac{IEMA}. Para mitigar o mal condicionamento desse problema inverso, aplica-se regularização L2 (Ridge), sendo comparados dois cenários: calibração padrão e calibração com regularização.

Na segunda etapa, as emissões sintéticas horárias geradas pela equação calibrada tornam-se o alvo (\textit{target}) para o treinamento de uma rede neural \ac{MLP}. A rede aprende diretamente a relação entre as variáveis operacionais (geração, categoria de operação, fase de transição), as características estáticas da usina (combustível, potência instalada, eficiência) e a emissão horária, sem incorporar explicitamente a equação física na função de perda.

Essa abordagem em cascata permite: (i) garantir consistência com os totais anuais do \ac{IEMA}; (ii) capturar não-linearidades que a equação paramétrica pode não expressar plenamente; e (iii) avaliar, por meio da comparação entre os cenários padrão e L2, qual estratégia de calibração produz alvos sintéticos mais coerentes para o aprendizado supervisionado.

\section{Fontes de Dados e Descrição das Variáveis}

Foram empregadas duas fontes principais de dados, abrangendo o período de 2020 a 2024.

\subsection{Operador Nacional do Sistema Elétrico (\ac{ONS})}

Por meio da iniciativa Dados Abertos, obtém-se a geração horária de todas as usinas do \ac{SIN}. Os registros incluem: identificação da usina, data, hora e geração horária em MW. O conjunto é filtrado para reter apenas usinas termelétricas \cite{ons_dados}.

\subsection{Instituto de Energia e Meio Ambiente (\ac{IEMA})}

Os inventários anuais do \ac{IEMA} fornecem \cite{iema2022, iema2023, iema2024}:
\begin{itemize}
    \item Emissões totais anuais de \ac{tCO$_2$e} por usina;
    \item Características cadastrais fixas: potência instalada (MW), fator de capacidade (\%), eficiência energética (\%), tipo de combustível e ciclo de operação (ciclo combinado, ciclo simples, cogeração).
\end{itemize}

Usinas de cogeração são excluídas da análise por não apresentarem, nos inventários do \ac{IEMA}, informações completas de potência instalada e eficiência — características indispensáveis à modelagem. Após a integração das fontes e a remoção de registros faltantes, restam 39 usinas termelétricas com dados completos para o período 2020-2024. As emissões anuais são mantidas em \ac{tCO$_2$e}, e a geração horária em MW.

\section{Tratamento das Sazonalidades Temporais}

As variáveis temporais \textit{Índice} (posição da hora dentro do ano) e \textit{Ano} não podem ser tratadas como grandezas ordinais lineares, pois a relação entre instantes próximos ao final e ao início de um ano ou período não é capturada por uma codificação escalar convencional.

Para preservar a periodicidade natural do calendário, adota-se uma codificação polar (seno/cosseno). Dado um índice \(i \in [1, H]\), onde \(H \in \{8760, 8784\}\) representa as horas do respectivo ano (considerando anos bissextos), definem-se as seguintes transformações:

\begin{equation}
\text{Seno\_Índice}(i) = \sin\left(\frac{2\pi i}{H}\right), \qquad
\text{Cosseno\_Índice}(i) = \cos\left(\frac{2\pi i}{H}\right)
\label{eq:codificacao_indice}
\end{equation}

Analogamente, para o ano \(a \in [2020, 2024]\), define-se:

\begin{equation}
\text{Seno\_Ano}(a) = \sin\left(\frac{2\pi (a - 2020)}{4}\right), \qquad
\text{Cosseno\_Ano}(a) = \cos\left(\frac{2\pi (a - 2020)}{4}\right)
\label{eq:codificacao_ano}
\end{equation}

Essa representação mapeia valores temporais para pontos sobre o círculo unitário, de modo que instantes próximos (ex: hora 8760 e hora 1) tornam-se vizinhos no espaço polar, refletindo sua continuidade temporal.

\section{Modelagem das Transições Operacionais}

Termelétricas operam tipicamente em patamares discretos de carga, não em regime continuamente variável. Para capturar esse comportamento, define-se um sistema de categorias de geração com base na proximidade relativa entre valores de potência.

\subsection{Definição de Categorias}

Dada a série de geração horária \(G = [g_1, g_2, \ldots, g_T]\) em MW, valores inferiores a 20 MW são classificados na categoria 0 (carga muito baixa ou desligamento). Para os demais valores, dois pontos pertencem à mesma categoria se sua diferença relativa for inferior ou igual a um limiar \(\theta = 15\%\) do maior valor do grupo.

Formalmente, \(g_i\) e \(g_j\) pertencem à mesma categoria se:

\begin{equation}
\frac{|g_i - g_j|}{\max(g_i, g_j)} \leq \theta
\label{eq:diferenca_relativa}
\end{equation}

O limiar de 15\% foi definido com base na análise exploratória da distribuição dos incrementos de carga nos dados do \ac{ONS}, representando o valor típico que separa regimes operacionais distintos sem fragmentar artificialmente a operação.

\subsection{Identificação de Transições}

Seja \(c_t\) a categoria no instante \(t\). Define-se a moda \(m_t\) das categorias em uma janela de três horas anteriores e três horas posteriores, excluindo o próprio instante:

\begin{equation}
m_t = \text{moda}\{c_{t-3}, c_{t-2}, c_{t-1}, c_{t+1}, c_{t+2}, c_{t+3}\}
\label{eq:transicao}
\end{equation}

O instante \(t\) é classificado como \textbf{transição} se \(c_t \neq m_t\). Este critério captura períodos de rampa de carga (partidas, paradas e variações controladas), nos quais o fator de emissão horária pode diferir substancialmente do regime estabilizado \cite{kumar2012}.

\subsection{Duração e Fase da Transição}

Para cada evento de transição contínuo (sequência de horas consecutivas classificadas como transição), definem-se:
\begin{itemize}
    \item \textbf{Duração} \(d\): número de horas consecutivas do evento;
    \item \textbf{Fase} \(f\): posição relativa dentro do evento, classificada como 1 (início), 2 (meio) ou 3 (fim). Transições com duração \(d = 1\) recebem \(f = 1\) (início e fim coincidentes).
\end{itemize}

Essas variáveis permitem que o modelo compreenda a transição como um processo contínuo, com dinâmica de emissão potencialmente distinta ao longo de suas fases.

\section{Formulação do Problema de Calibração}

A relação entre geração horária e emissão horária é modelada por uma equação com termo quadrático e termo exponencial, conforme proposta na literatura \cite{gent1971minimum, talaq1994summary}:

\begin{equation}
E_k^{\text{horária}}(P_{G_i}) = \alpha_k P_{G_i}^2 + \beta_k P_{G_i} + \gamma_k + \omega_k e^{\mu_k P_{G_i}}
\label{eq:emissao_horaria}
\end{equation}

onde:
\begin{itemize}
    \item \(P_{G_i}\) representa a geração em MW, dividida por 100 para normalização dos dados, garantindo estabilidade numérica no termo exponencial;
    \item \(\alpha_k, \beta_k, \gamma_k, \omega_k, \mu_k\) são coeficientes a serem calibrados por usina;
    \item O subscrito \(k\) indica a usina.
\end{itemize}

\subsection{Calibração Padrão}

Os coeficientes são obtidos resolvendo o seguinte problema de minimização, que busca minimizar o erro quadrático médio entre a emissão anual estimada (somatório das emissões horárias) e a emissão anual real fornecida pelo \ac{IEMA}:

\begin{equation}
\min_{\alpha,\beta,\gamma,\omega,\mu} \frac{1}{N} \sum_{j=1}^{N} \left( \sum_{i=1}^{H} E_{\text{usina}}^{i}(P_{G_i}) - E_{\text{usina}}^{j} \right)^2
\label{eq:minimizacao_padrao}
\end{equation}

sujeito às restrições físicas:

\begin{equation}
\alpha \geq 0, \quad \mu \leq \mu_{\max}
\label{eq:restricoes_fisicas}
\end{equation}

onde:
\begin{itemize}
    \item \(N\) é o número de usinas (\(N = 39\));
    \item \(H\) é o número de horas do respectivo ano (\(H = 8760\) para anos não bissextos, \(H = 8784\) para anos bissextos);
    \item \(E_{\text{usina}}^{j}\) é a emissão anual real do \ac{IEMA} para a usina \(j\);
    \item \(\mu_{\max}\) é um limite superior para evitar \textit{overflow} numérico do termo exponencial (definido empiricamente em \(0,5\)).
\end{itemize}

Apenas o coeficiente \(\alpha\) foi explicitamente restringido à não-negatividade, por representar o termo quadrático associado ao crescimento da emissão em cargas elevadas. Os coeficientes \(\gamma\) e \(\omega\) foram mantidos livres para permitir maior flexibilidade durante a calibração do problema inverso \cite{wood2013power}. O coeficiente \(\mu\) recebe um limite superior \(\mu_{\max} = 0,5\) para evitar crescimento excessivo do termo exponencial, que poderia causar instabilidade numérica.

\subsection{Regularização L2 (Ridge)}

O problema de calibração inversa é intrinsecamente mal condicionado, podendo produzir coeficientes excessivamente elevados que se cancelam mutuamente (\textit{overfitting} na calibração). Para mitigar esse efeito, introduz-se um termo de regularização L2 (Tikhonov) à função objetivo \cite{tikhonov1977solutions}:

\begin{equation}
\min_{\alpha,\beta,\gamma,\omega,\mu} \frac{1}{N} \sum_{j=1}^{N} \left( \sum_{i=1}^{H} E_{\text{usina}}^{i}(P_{G_i}) - E_{\text{usina}}^{j} \right)^2 + \lambda \left( \alpha^2 + \beta^2 + \gamma^2 + \omega^2 + \mu^2 \right)
\label{eq:regularizacao_l2}
\end{equation}

sujeito às mesmas restrições da calibração padrão (Equação \ref{eq:restricoes_fisicas}).

O parâmetro \(\lambda \geq 0\) controla a intensidade da regularização: \(\lambda = 0\) recupera o problema padrão (Equação \ref{eq:minimizacao_padrao}), enquanto valores positivos penalizam coeficientes de grande magnitude, promovendo soluções mais estáveis e generalizáveis. O valor de \(\lambda\) deve ser selecionado de forma a equilibrar a redução do erro de calibração e a estabilidade dos coeficientes. Neste trabalho, adotou-se \(\lambda = 30\), conforme justificado na Seção 4.3.3.

\section{Emissões Sintéticas como Alvos para Aprendizado}

Uma vez calibrados os coeficientes (para os cenários padrão e com regularização L2), aplica-se a Equação \ref{eq:emissao_horaria} sobre a série horária de geração para produzir as \textbf{emissões sintéticas horárias} \(\hat{E}_t\). Essas emissões sintéticas tornam-se a variável alvo para o treinamento supervisionado da rede neural.

É importante observar que a rede neural, ao ser treinada com esses alvos sintéticos, herda os erros de modelagem da equação paramétrica e da calibração. Conforme discutido na literatura \cite{nowok2016synthpop, stenger2024synthetic}, os dados sintéticos gerados por modelos paramétricos refletem integralmente as premissas e limitações do modelo utilizado. Portanto, a comparação entre os cenários padrão e L2 não avalia apenas a qualidade da rede, mas também qual estratégia de calibração produz alvos mais consistentes para o aprendizado supervisionado.

A fidelidade da base sintética depende diretamente da qualidade do ajuste e da capacidade do modelo paramétrico em representar o comportamento real de emissão da usina.

\section{Arquitetura da Rede Neural}

Optou-se por uma arquitetura do tipo \ac{MLP} em detrimento de redes recorrentes (\ac{LSTM}, \ac{GRU}) porque a emissão horária, no modelo proposto, depende essencialmente da geração no mesmo instante \(t\), e não de uma sequência temporal longa. A \ac{MLP} é suficiente, mais simples e menos suscetível a \textit{overfitting} com o volume de dados disponível \cite{goodfellow2016deep}.


\subsection{Modelos Individualizados por Usina}

Um aspecto fundamental da metodologia é que um modelo de rede neural \ac{MLP} separado é treinado para cada uma das 39 usinas. Esta abordagem individualizada é necessária porque cada usina possui características operacionais próprias: diferentes tipos de combustível, eficiências energéticas, potências instaladas e perfis de transição entre patamares de carga. Um modelo único para todas as usinas não seria capaz capturar essas particularidades. Portanto, o processo descrito na Seção 3.7.1 é repetido individualmente para cada usina, resultando em 39 modelos de rede neural distintos ao final do treinamento.

\subsection{Estrutura das Camadas}

A Tabela \ref{tab:arquitetura} apresenta a arquitetura adotada.

\begin{table}[H]
    \centering
    \renewcommand{\arraystretch}{1.5}
    \begin{tabular}{|c|c|c|c|}
        \hline
        \textbf{Camada} & \textbf{Neurônios} & \textbf{Ativação} & \textbf{Regularização} \\
        \hline
        Entrada & \(d_{\text{in}}\) & -- & -- \\
        Oculta 1 & 128 & \ac{ReLU} & Dropout (0,3) \\
        Oculta 2 & 64 & \ac{ReLU} & Dropout (0,3) \\
        Oculta 3 & 32 & \ac{ReLU} & Dropout (0,3) \\
        Saída & 1 & \ac{ReLU} & -- \\
        \hline
    \end{tabular}
    \caption{Arquitetura da Rede Neural MLP}
    \label{tab:arquitetura}
\end{table}

A arquitetura em funil (128 $\rightarrow$ 64 $\rightarrow$ 32) reduz progressivamente a capacidade representacional, atuando como uma forma implícita de regularização e prevenindo \textit{overfitting}. A função de ativação \ac{ReLU} é escolhida por sua simplicidade computacional e por mitigar o problema do gradiente evanescente em comparação com funções sigmoidais \cite{goodfellow2016deep}. O \textit{dropout} com taxa de \(0,3\) desativa aleatoriamente \(30\%\) dos neurônios a cada iteração de treinamento, forçando a rede a aprender representações redundantes e robustas \cite{srivastava2014dropout}.

\subsection{Estratégias de Regularização Adicionais}

Além do \textit{dropout}, empregam-se as seguintes estratégias de regularização:

\begin{itemize}
    \item \textbf{Early stopping}: o treinamento é interrompido quando a perda no conjunto de validação não apresenta melhora por um número pré-definido de épocas (paciência), restaurando os pesos da melhor época. Isso evita \textit{overfitting} ao interromper o treinamento antes que a rede memorize ruídos \cite{goodfellow2016deep}.

    \item \textbf{Gradient clipping}: a norma L2 do vetor de gradientes é limitada a um valor máximo antes da atualização dos pesos. Esta técnica previne a explosão dos gradientes, fenômeno que pode ocorrer quando a função de perda apresenta regiões de alta curvatura ou quando amostras com pesos elevados produzem gradientes desproporcionais \cite{goodfellow2016deep}.
\end{itemize}

\section{Variáveis de Entrada e Pré-processamento Conceitual}

As variáveis de entrada da rede neural são agrupadas em duas categorias: estáticas e dinâmicas.

\subsection{Variáveis Estáticas}

Variáveis estáticas são invariantes no tempo para cada usina:
\begin{itemize}
    \item Tipo de combustível (categórica, sem ordem intrínseca);
    \item Ciclo de operação (categórica);
    \item Potência instalada (MW) — grandeza com significado físico direto;
    \item Eficiência energética (\%) — grandeza com significado físico direto;
    \item Fator de capacidade (\%) — grandeza com significado físico direto.
\end{itemize}

\subsection{Variáveis Dinâmicas}

Variáveis dinâmicas variam a cada hora:
\begin{itemize}
    \item Geração horária (MW);
    \item Categoria de geração (categórica, proveniente da classificação descrita na Seção 4);
    \item Duração da transição (h);
    \item Fase da transição (1, 2 ou 3);
    \item Índice temporal (codificado polarmente conforme Seção 3);
    \item Ano (codificado polarmente conforme Seção 3).
\end{itemize}

\subsection{Tratamento das Variáveis}

Variáveis categóricas (combustível, ciclo de operação, categoria de geração) são transformadas por \textit{one-hot encoding}, que as converte em vetores binários sem impor ordenação artificial.

Variáveis contínuas de diferentes ordens de grandeza (geração, duração, fase) requerem padronização para que todas contribuam igualmente durante o treinamento. Adota-se a padronização \textit{z-score}:

\begin{equation}
x_{\text{pad}} = \frac{x - \mu}{\sigma}
\label{eq:padronizacao}
\end{equation}

onde \(\mu\) e \(\sigma\) são a média e o desvio padrão estimados \textbf{exclusivamente} nos dados de treinamento, para evitar vazamento de informações do futuro.

\section{Métricas de Avaliação}

Para avaliar a qualidade das estimativas, empregam-se as seguintes métricas \cite{chai2014root, willmott2005advantages, myttenaere2016mean}.

\subsection{Erro Absoluto Médio (\ac{MAE})}

\begin{equation}
\text{MAE} = \frac{1}{n} \sum_{i=1}^{n} |y_i - \hat{y}_i|
\label{eq:mae}
\end{equation}
Interpretável diretamente na unidade original (tCO\(_2\)/h).

\subsection{Erro Quadrático Médio (\ac{MSE})}

\begin{equation}
\text{MSE} = \frac{1}{n} \sum_{i=1}^{n} (y_i - \hat{y}_i)^2
\label{eq:mse}
\end{equation}
Penaliza mais severamente erros grandes, sendo útil para detectar a presença de \textit{outliers}.

\subsection{Raiz do Erro Quadrático Médio (\ac{RMSE})}

\begin{equation}
\text{RMSE} = \sqrt{\text{MSE}}
\label{eq:rmse}
\end{equation}
Devolve a métrica à unidade original (tCO\(_2\)/h), combinando a sensibilidade a \textit{outliers} do \ac{MSE} com a interpretabilidade do \ac{MAE}.

\subsection{Razão de Estabilidade}

A razão $R = \text{MSE} / \text{MAE}^2$ é empregada como indicador da estabilidade do modelo. Esta razão está diretamente relacionada à distribuição dos erros e à presença de \textit{outliers}. Para uma distribuição normal de erros, a razão converge para $\pi/2 \approx 1,57$, pois $\text{MSE} = \sigma^2$ e $\text{MAE} = \sigma \sqrt{2/\pi}$, resultando em $\text{MSE}/\text{MAE}^2 = \pi/2$.

Com base nesta propriedade e em considerações práticas, adotam-se os seguintes limiares:

\begin{itemize}
    \item $R < 3$: modelo estável, erros concentrados em torno da média (até aproximadamente o dobro do valor teórico);
    \item $3 \leq R \leq 5$: instabilidade moderada, presença de alguns \textit{outliers};
    \item $R > 5$: presença significativa de \textit{outliers}, modelo pouco confiável.
\end{itemize}

Estes limiares foram definidos com base na análise exploratória dos dados e na experiência prática do autor, sendo utilizados como referência para comparação entre os cenários avaliados.